\section{Conclusiones}

El modelo resuelve el problema planteado. En las diez instancias de prueba, el óptimo que demuestra Gecode coincide con el de HiGHS y, en las tres más pequeñas, con el cálculo manual, y la instancia infactible se detecta sin ramificar; además, el solucionador independiente de Detalles de implementación coincidió con las cinco implementaciones en 6\,000 corridas. El elemento decisivo para Gecode son las restricciones de dominancia: sin ellas no demuestra el óptimo de ninguna instancia de tamaño real en 60~s, y con ellas el ejemplo del enunciado se resuelve en 9 nodos. La ruptura de simetrías reduce el árbol solo cuando hay centrales o clientes idénticos, y las restricciones redundantes, aunque correctas, no podan, porque lo que deducen ya lo deducen R4, R5 y la dominancia. Por eso se recomienda \texttt{impl} 5: incluye la dominancia y la ruptura de simetrías, y las redundantes solo agregan algo de tiempo por nodo.

Las mediciones muestran que la anotación de búsqueda influye tanto en la calidad de las soluciones encontradas como en la posibilidad de demostrar que son óptimas. La estrategia por orden de mérito fue la más consistente: encontró el óptimo de \texttt{exp\_01} y \texttt{exp\_08}, fue la única que demostró el de \texttt{PlantaEnergia\_08} y dio las mejores soluciones de Gecode en las instancias grandes. No obstante, las corridas de 14, 25 y 30 días terminaron sin prueba de optimalidad, y \texttt{dom\_w\_deg} con división de dominios no encontró soluciones en las instancias más grandes dentro de 60~s, aunque exploró menos nodos que la estrategia 2 en cuatro instancias pequeñas. Por esto, los resultados no justifican afirmar que una estrategia sea superior para toda instancia.

HiGHS y COIN-BC demostraron los mismos óptimos en las doce instancias experimentales y permitieron cuantificar las brechas de Gecode en las grandes: 0,46~\% y 0,30~\%. Esto respalda el uso de solucionadores MIP como referencia exacta para este modelo. También confirma que \texttt{sat} solo indica que se encontró una solución factible; no debe interpretarse como una prueba de optimalidad.

En conjunto, las instancias pequeñas sirven para verificar el comportamiento del modelo y comparar las anotaciones, mientras que las grandes exhiben el límite práctico de Gecode bajo el plazo adoptado. La estrategia 2 es una elección razonable cuando se usa Gecode y se busca una buena solución dentro de un tiempo acotado; para certificar el óptimo en las instancias grandes probadas, HiGHS fue la alternativa más rápida. Estas conclusiones se restringen a los datos, implementaciones y límite de 60~s de los experimentos realizados.
